% ---
% id: reductions/prop:common
% kind: prop
% title: "Common and dominating targets"
% summary: >-
%   If every branch target squashes to one POVM, the coarse-grained measurement
%   of [[lem:master]] squashes to that POVM directly, through the sum of the
%   composed branch channels. The branch blocks then disappear from the
%   target.
% prominence: standard
% status: accepted
% uses: [lem:master]
% mentions: [thm:mixture, cor:register-instance, cor:instances]
% assumes: []
% refs: []
% provenance: newly-derived
% lean:
% ---

Only in Corollary~\ref{cor:instances}\,(i) is the branch weight state
independent, hence the constraint a fixed number. In its case~(ii) the weight
is the projection onto the block the input fell in, and that operator-valued
case is the one a flag-state squasher runs on. Its case~(iii) is stated for a
general POVM $\{N_j\}$ rather than a PVM, and it is what the
register-conditioned squashing of Theorem~\ref{thm:mixture} runs on: a
\emph{classical branch register} read by a POVM rather than a PVM, which
Corollary~\ref{cor:register-instance} lifts to \tref[reduction instance]{reduction instances}. The price
of the generality is that each branch gets its own block of $\Qsys$ rather
than sharing one register, which is exactly the block structure of $\Qsys$
the consuming declaration wants anyway. Where every branch target squashes to
one common target, the blocks can be merged again.

\begin{proposition}\label{prop:common}
In the setting of Lemma~\ref{lem:master}, suppose additionally
$F^{(j)}\dsq_{\Psi_j}F$ for a single POVM $F$ on a space $\Qsys_{*}$ (a
\emph{dominating target}; $\Psi_j=\idmap$ if all branch targets already
coincide). Then
\begin{equation*}
  \Big\{\sum_{j}M_{v,j}\Big\}_v\;\dsq_{\Lambda'}\;F,
  \qquad
  \Lambda'\coloneqq\sum_{j\in J}\Psi_j\circ\Lambda_j\circ\Ins_j .
\end{equation*}
\end{proposition}

\begin{proof}
$\Lambda'$ is CP as a sum of compositions of CP maps and trace preserving by
the computation in Lemma~\ref{lem:master}\,(ii). Then
$\Lambda'^{\dagger}(F_v)=\sum_j\Ins_j^{\dagger}\Lambda_j^{\dagger}
\Psi_j^{\dagger}(F_v)
=\sum_j\Ins_j^{\dagger}\Lambda_j^{\dagger}\big(F^{(j)}_v\big)
=\sum_j\Ins_j^{\dagger}\big(M^{(j)}_v\big)=\sum_jM_{v,j}$.
\end{proof}
